\section{Conclusion}
\label{sec:conclusion}
This work examined one specific claim of the Adios paper---that yielding during a remote page fetch outperforms spinning once context-switch cost is small---by constructing a discrete-event simulator and validating it against closed-form capacity bounds. With a 20\% local hit rate, our experiments showed that yielding lifts peak throughput from 1.77 to 2.50~MRPS (a factor of 1.41) and roughly halves the 99.9th percentile latency at 1.5~MRPS. The fundamental reason in both cases is that the yield system's ceiling is the network link, while the busy-wait system's ceiling is its processor pool---the same explanation the original paper offers. Yielding loses its advantage when two context switches together cost as much as one fetch, which in our configuration happens at $t_s = 1.5~\mu$s, roughly 75 times larger than the switch cost of a unithread design. Using the capacity formulas as a predictive model rather than just a check, we find that yielding needs fewer than half as many workers as spinning to carry 2.5~MRPS (5.1 versus 11.0) and that its advantage grows with fetch time---a prediction we did not test with the simulator. The page-fault-aware dispatch rule helped by at most 1.2\% in our model, far below the gains the original system reports, and our memory-shrink experiment did not reproduce the published throughput-loss pattern.

Going beyond the reproduction exercise, we identified a fundamental weakness of pure yielding: it wastes context-switch budget when the system is lightly loaded and becomes sensitive to rising switch costs in environments without unithread-style scheduling. To address this, we introduced \emph{Load-Aware Adaptive Hybrid Yielding}, which dynamically switches between spin-polling when uncontended and yielding when queue pressure builds. The result is the best of both worlds: the median latency of spinning (5.30~$\mu$s) at light load, the saturation throughput (2.50~MRPS) and tail protection of yielding at heavy load, and complete immunity to context-switch cost increases up to 1.5~$\mu$s.

The simulator does not reproduce Adios---it omits the unikernel, unithread layout, polling delegation, reclaimer, and real network hardware, using fixed service times throughout. Its results are therefore best understood as evidence about the scheduling mechanism itself and as a convenient tool for identifying break-even operating points, rather than as predictions of absolute performance on real hardware. All code, result files, and figures are available in the project repository so that every number in this report can be independently regenerated.
